\documentclass[tikz,border=0pt]{standalone}
\usepackage{ravel-figures}
\begin{document}
\begin{tikzpicture}[x=1mm,y=1mm]
\tikzset{b/.style={box=#1, minimum height=16mm}, b/.default=48mm}
\ravelheader{The physics toolchain}{Three chains of established tools, run with versions recorded.}
\tikzset{t/.style={box=37mm, minimum height=18mm, draw=rvhepline, fill=rvhepfill, anchor=north west},
         i/.style={box=37mm, minimum height=18mm, draw=rvpeopleline, fill=rvpeoplefill, anchor=north west},
         lane/.style={anchor=north west, inner sep=0pt, font=\ravelSmallFont\bfseries, text=rvmuted}}
\node[lane] at (5,-23) {Collider};
\node[i] (c0) at (5,-28) {\nodetitle{Model}\nodesub{cards or model files}};
\node[t] (c1) at (46,-28) {\nodetitle{Events}\nodesub{MadGraph, Pythia}};
\node[t] (c2) at (87,-28) {\nodetitle{Detector}\nodesub{Delphes, FastJet}};
\node[t] (c3) at (128,-28) {\nodetitle{Analysis}\nodesub{Rivet, CheckMATE, SModelS}};
\node[lane] at (5,-52) {Theory};
\node[i] (t0) at (5,-57) {\nodetitle{Lagrangian}\nodesub{your model}};
\node[t] (t1) at (46,-57) {\nodetitle{Model files}\nodesub{FeynRules, SARAH}};
\node[t] (t2) at (87,-57) {\nodetitle{Spectrum}\nodesub{FlexibleSUSY, SPheno}};
\node[t] (t3) at (128,-57) {\nodetitle{Diagrams}\nodesub{FeynArts, FeynCalc}};
\node[lane] at (5,-81) {Cosmology};
\node[i] (k0) at (5,-86) {\nodetitle{Parameters}\nodesub{model and data}};
\node[t] (k1) at (46,-86) {\nodetitle{Solvers}\nodesub{CAMB, CLASS}};
\node[t] (k2) at (87,-86) {\nodetitle{Inference}\nodesub{Cobaya}};
\node[t] (k3) at (128,-86) {\nodetitle{Simulation}\nodesub{Nyx on a cluster}};
\foreach \p in {c,t,k} { \draw[flow] (\p0) -- (\p1); \draw[flow] (\p1) -- (\p2); \draw[flow] (\p2) -- (\p3); }
\ravellegend{-112}{people/Your inputs, hep/Physics tools}
\ravelcanvasfit
\end{tikzpicture}
\end{document}
